\SourceLeaf{220}%
\noindent oder wenn man lieber will, der symmetrischen
\[
EBC'+GCA'+HAB'-ECB'-GAC'-HBA'
\]
\noindent abhangen.

Die Vergleichung der Darstellungen in 3 und 4 beruht auf ganz ähnlichen Gründen, wie die von 2 und 3, und die ähnliche oder verkehrte Lage der Theile hängt von dem positiven oder negativen Zeichen der Grösse
\[
(\frac{\mathrm{d}p}{\mathrm{d}t})\mathord{.}(\frac{\mathrm{d}q}{\mathrm{d}u})-(\frac{\mathrm{d}p}{\mathrm{d}u})\mathord{.}(\frac{\mathrm{d}q}{\mathrm{d}t})
\]
\noindent ab; und eben so bestimmt das positive oder negative Zeichen von
\[
(\frac{\mathrm{d}P}{\mathrm{d}T})\mathord{.}(\frac{\mathrm{d}Q}{\mathrm{d}U})-(\frac{\mathrm{d}P}{\mathrm{d}U})\mathord{.}(\frac{\mathrm{d}Q}{\mathrm{d}T})
\]
\noindent die ähnliche oder verkehrte Lage der Theile in den Darstellungen 5 und 6.

Was endlich die Vergleichung der Darstellungen 4 und 5 unter sich betrifft, so können wir uns auf die Analyse des 8. Artikels beziehen, aus welcher erhellet, dass jene in den kleinsten Theilen ähnlich, oder verkehrt liegend sind, je nachdem man die erste oder zweite Auflösung gewählt, d. i. entweder
\[
P+iQ=f(p+iq)\quad\text{und}\quad P-iQ=f'(p-iq)
\]
\noindent oder
\[
P+iQ=f(p-iq)\quad\text{und}\quad P-iQ=f'(p+iq)
\]
\noindent gesetzt hat.

Aus diesem allen ziehen wir nunmehro den Schluss, dass man, wenn die Darstellung auf der Fläche, deren Gleichung \(\Psi=0\) ist, dem Urbilde auf der Fläche, deren Gleichung \(\psi=0\) ist, in den kleinsten Theilen nicht bloss ähnlich, sondern auch ähnlich liegend sein soll, auf die Anzahl der negativen Grössen, welche unter diesen vier Grössen vorkommen,
\[
\begin{gathered}
\frac{ab'-ba'}{h},\quad (\frac{\mathrm{d}p}{\mathrm{d}t})\mathord{.}(\frac{\mathrm{d}q}{\mathrm{d}u})-(\frac{\mathrm{d}p}{\mathrm{d}u})\mathord{.}(\frac{\mathrm{d}q}{\mathrm{d}t}),\\
(\frac{\mathrm{d}P}{\mathrm{d}T})\mathord{.}(\frac{\mathrm{d}Q}{\mathrm{d}U})-(\frac{\mathrm{d}P}{\mathrm{d}U})\mathord{.}(\frac{\mathrm{d}Q}{\mathrm{d}T}),\quad \frac{AB'-BA'}{H}
\end{gathered}
\]
\noindent Rücksicht nehmen muss; ist gar keine oder eine gerade Anzahl darunter, so wird die erste; ist eine oder drei negative unter ihnen, so wird die zweite Auflösung gewählt werden müssen. Bei entgegengesetzter Wahl findet allemal eine verkehrte Aehnlichkeit Statt.
